% =====================================================================
\chapter{Taylor, hessiano y puntos críticos}
\label{tema:4}
% =====================================================================

El gradiente describe la inclinación de una función; la curvatura la describe el hessiano. Este tema introduce las derivadas de segundo orden, demuestra el desarrollo de Taylor de segundo orden en varias variables, lo usa para clasificar puntos críticos y, como aplicación central, determina para qué tasas de aprendizaje converge el descenso de gradiente sobre una función cuadrática. Se sigue~\cite[§5.7 y §5.8]{mml}, \cite[§4.2 y §4.3.1]{gbc} y la clase~9 de~\cite{mit1802}.

\section{El hessiano}

\begin{definicion}[Hessiano]\label{def:hessiano}
Sea $f\colon U\subseteq\R^{n}\to\R$ con derivadas parciales segundas. El \emph{hessiano} de $f$ en $x$ es la matriz $H_f(x)=\nabla^{2}f(x)\in\R^{n\times n}$ de entradas
\[
\bigl(H_f(x)\bigr)_{ij}=\frac{\partial^{2}f}{\partial x_i\,\partial x_j}(x).
\]
Equivalentemente, $H_f$ es el jacobiano del gradiente: $H_f=J_{\nabla f}$.
\end{definicion}

\begin{teorema}[Schwarz]\label{teo:schwarz}
Si $f\in C^{2}(U)$, entonces $\partial^{2}f/\partial x_i\partial x_j=\partial^{2}f/\partial x_j\partial x_i$ en todo punto de $U$. En consecuencia, $H_f(x)$ es simétrica.
\end{teorema}

La demostración se encuentra en los textos de análisis de varias variables y se omite. Al ser simétrico, el hessiano es diagonalizable en una base ortonormal con autovalores reales (teorema espectral, módulo~3), y esos autovalores son las curvaturas de $f$ en las direcciones de los autovectores~\cite[§4.3.1]{gbc}.

\begin{ejemplo}\label{ej:hessiano}
Para $f(x,y)=x^{2}y+3y$ se tiene $\partial^{2}f/\partial x^{2}=2y$, $\partial^{2}f/\partial x\partial y=2x$ y $\partial^{2}f/\partial y^{2}=0$, luego
\[
H_f(2,1)=\begin{pmatrix}2&4\\4&0\end{pmatrix},
\]
con autovalores $1\pm\sqrt{17}$, es decir, $5{,}123$ y $-3{,}123$.
\end{ejemplo}

\begin{observacion}[Tamaño del hessiano]
Para $N$ parámetros el hessiano tiene $N^{2}$ entradas. Con $N=7\cdot 10^{9}$ serían $4{,}9\cdot 10^{19}$ números, algo imposible de almacenar. En la práctica solo se calculan productos hessiano-vector $H v$, que cuestan un múltiplo pequeño de un gradiente, o aproximaciones diagonales.
\end{observacion}

\section{Desarrollo de Taylor de segundo orden}

\begin{teorema}[Taylor de segundo orden]\label{teo:taylor}
Sea $f\in C^{2}(U)$ y sean $x\in U$ y $h\in\R^{n}$ tales que el segmento $\{x+th:t\in[0,1]\}$ está contenido en $U$. Entonces:
\begin{enumerate}
  \item (forma de Lagrange) existe $\tau\in(0,1)$ tal que
  \[
  f(x+h)=f(x)+\nabla f(x)\T h+\tfrac12\,h\T H_f(x+\tau h)\,h;
  \]
  \item (forma local) cuando $h\to 0$,
  \[
  f(x+h)=f(x)+\nabla f(x)\T h+\tfrac12\,h\T H_f(x)\,h+o(\norm{h}^{2}).
  \]
\end{enumerate}
\end{teorema}

\begin{proof}
Sea $\varphi(t)=f(x+th)$ para $t\in[0,1]$. Por el \cref{lem:curva}, $\varphi'(t)=\nabla f(x+th)\T h=\sum_i h_i\,\partial_i f(x+th)$. Aplicando de nuevo el lema a cada función $\partial_i f$, que es diferenciable por ser $f\in C^{2}$,
\[
\varphi''(t)=\sum_{i}h_i\sum_{j}h_j\,\partial_j\partial_i f(x+th)=h\T H_f(x+th)\,h.
\]
El teorema de Taylor de una variable con resto de Lagrange da $\varphi(1)=\varphi(0)+\varphi'(0)+\tfrac12\varphi''(\tau)$ para algún $\tau\in(0,1)$, que es (i). Para (ii), se resta el término $\tfrac12 h\T H_f(x)h$:
\[
\Bigl|\tfrac12\,h\T\bigl(H_f(x+\tau h)-H_f(x)\bigr)h\Bigr|\le\tfrac12\,\norm{H_f(x+\tau h)-H_f(x)}\,\norm{h}^{2},
\]
y la norma de la diferencia tiende a~0 cuando $h\to 0$ por la continuidad de las segundas derivadas.
\end{proof}

\begin{ejemplo}\label{ej:taylor}
Con la función del \cref{ej:hessiano} en $x=(2,1)$ y $h=(0{,}1;\,-0{,}1)$, donde $f=7$ y $\nabla f=(4,7)\T$:
\begin{enumerate}[label=\arabic*.]
  \item primer orden: $7+(4\cdot 0{,}1-7\cdot 0{,}1)=6{,}7$;
  \item término cuadrático: $h\T Hh=2\cdot 0{,}01+2\cdot 4\cdot(0{,}1)(-0{,}1)+0=-0{,}06$, luego el segundo orden da $6{,}7-0{,}03=6{,}67$;
  \item valor exacto: $f(2{,}1;\,0{,}9)=4{,}41\cdot 0{,}9+2{,}7=6{,}669$.
\end{enumerate}
El error baja de $0{,}031$ a $0{,}001$ al incluir la curvatura.
\end{ejemplo}

\section{Clasificación de puntos críticos}

\begin{definicion}[Punto crítico]
Un punto $x^{*}\in U$ es \emph{crítico} (o estacionario) si $\nabla f(x^{*})=0$. Es un \emph{mínimo local} si existe $r>0$ con $f(x)\ge f(x^{*})$ para $\norm{x-x^{*}}<r$, y un \emph{mínimo local estricto} si la desigualdad es estricta para $x\ne x^{*}$. Un punto crítico que no es mínimo ni máximo local es un \emph{punto de silla}.
\end{definicion}

\begin{teorema}[Condición necesaria de primer orden]\label{teo:fermat}
Si $f$ es diferenciable en $x^{*}$ y $x^{*}$ es un mínimo local, entonces $\nabla f(x^{*})=0$.
\end{teorema}

\begin{proof}
Si $g=\nabla f(x^{*})\ne 0$, la \cref{prop:descenso} da puntos $x^{*}-\eta g$, arbitrariamente próximos a $x^{*}$, con $f(x^{*}-\eta g)<f(x^{*})$, lo que contradice que $x^{*}$ sea un mínimo local.
\end{proof}

\begin{teorema}[Criterio del hessiano]\label{teo:criterio}
Sea $f\in C^{2}(U)$ y $x^{*}$ un punto crítico.
\begin{enumerate}
  \item Si $H_f(x^{*})\succ 0$ (todos sus autovalores son positivos), $x^{*}$ es un mínimo local estricto.
  \item Si $H_f(x^{*})\prec 0$ (todos negativos), $x^{*}$ es un máximo local estricto.
  \item Si $H_f(x^{*})$ tiene autovalores de ambos signos, $x^{*}$ es un punto de silla.
\end{enumerate}
\end{teorema}

\begin{proof}
Sea $H=H_f(x^{*})$. Por el \cref{teo:taylor}(ii) con $\nabla f(x^{*})=0$, $f(x^{*}+h)-f(x^{*})=\tfrac12h\T Hh+o(\norm{h}^{2})$.

(i) Por el teorema espectral, $h\T Hh\ge\lambda_{\min}(H)\norm{h}^{2}$ con $\lambda_{\min}(H)>0$. Existe $r>0$ tal que el resto es menor que $\tfrac14\lambda_{\min}\norm{h}^{2}$ en valor absoluto si $0<\norm{h}<r$, y entonces $f(x^{*}+h)-f(x^{*})\ge\tfrac14\lambda_{\min}\norm{h}^{2}>0$.

(ii) Basta aplicar (i) a $-f$.

(iii) Sean $v_+$ y $v_-$ autovectores unitarios con autovalores $\lambda_+>0$ y $\lambda_-<0$. Para $h=t\,v_\pm$ se tiene $f(x^{*}+tv_\pm)-f(x^{*})=\tfrac12\lambda_\pm t^{2}+o(t^{2})$, que para $t\ne 0$ pequeño tiene el signo de $\lambda_\pm$. Así, en todo entorno de $x^{*}$ hay puntos con valores mayores y menores que $f(x^{*})$.
\end{proof}

\begin{table}[h]
\centering\small
\caption{Clasificación de un punto crítico según los autovalores de $H_f(x^{*})$.}\label{tab:criterio}
\begin{tabular}{@{}lll@{}}
\toprule
Autovalores de $H_f(x^{*})$ & Tipo de punto & Ejemplo en $(0,0)$ \\
\midrule
todos $>0$ & mínimo local estricto & $x^{2}+y^{2}$ \\
todos $<0$ & máximo local estricto & $-x^{2}-y^{2}$ \\
de ambos signos & punto de silla & $x^{2}-y^{2}$ \\
alguno nulo, el resto del mismo signo & no concluyente & $x^{4}+y^{2}$, $-x^{4}+y^{2}$ \\
\bottomrule
\end{tabular}
\end{table}

\begin{observacion}[Criterio en dos variables]
Para $n=2$, el producto de los autovalores es $\det H$ y su suma es $\tr H$. Por ello, si $\det H<0$ el punto es de silla, y si $\det H>0$ es un mínimo cuando $f_{xx}>0$ y un máximo cuando $f_{xx}<0$. Es la forma en que se presenta el criterio en~\cite[clase~9]{mit1802}.
\end{observacion}

\begin{observacion}[Conexión con los LLM]
El hessiano del \cref{ej:hessiano} tiene autovalores de ambos signos: es una curvatura de silla. En espacios de dimensión muy alta, la mayoría de los puntos críticos de la pérdida de una red son puntos de silla, no mínimos locales malos, porque basta con que uno solo de los $N$ autovalores sea negativo~\cite{dauphin}.
\end{observacion}

\section{Curvatura y descenso de gradiente}

El \cref{teo:taylor} permite analizar el descenso de gradiente más allá de la \cref{prop:descenso}. El caso modelo es una función cuadrática, que coincide con su propio desarrollo de Taylor de segundo orden.

\begin{proposicion}[Descenso de gradiente sobre una cuadrática]\label{prop:dg-cuadratica}
Sea $f(x)=\tfrac12\,x\T Ax$ con $A$ simétrica y autovalores $0<\lambda_1\le\dots\le\lambda_n$. La iteración $x_{k+1}=x_k-\eta\nabla f(x_k)$ converge al mínimo $x^{*}=0$ desde todo punto inicial si y solo si
\[
0<\eta<\frac{2}{\lambda_n}=\frac{2}{\lambda_{\max}(A)}.
\]
\end{proposicion}

\begin{proof}
Como $\nabla f(x)=Ax$, la iteración es $x_{k+1}=(I-\eta A)x_k$, luego $x_k=(I-\eta A)^{k}x_0$. Sea $A=Q\Lambda Q\T$ la descomposición espectral y $c^{(k)}=Q\T x_k$ las coordenadas en la base de autovectores. Entonces $c^{(k)}_i=(1-\eta\lambda_i)^{k}c^{(0)}_i$ para cada $i$. Todas las coordenadas tienden a~0 para todo $x_0$ si y solo si $\abs{1-\eta\lambda_i}<1$ para todo $i$, es decir, $0<\eta\lambda_i<2$. La condición más exigente es la del mayor autovalor.
\end{proof}

\begin{ejemplo}\label{ej:dg-cuad}
La función $f(x,y)=x^{2}+3y^{2}$ del \cref{ej:dg} es $\tfrac12x\T Ax$ con $A=\diag(2,6)$. El descenso converge si y solo si $\eta<2/6=1/3$. Con $\eta=0{,}1$, los factores de contracción son $1-0{,}2=0{,}8$ en $x$ y $1-0{,}6=0{,}4$ en $y$, lo que reproduce el paso $(1,1)\to(0{,}8;\,0{,}4)$.
\end{ejemplo}

\begin{observacion}[Número de condición]
El factor de contracción en la dirección $i$ es $\abs{1-\eta\lambda_i}$. Con la elección óptima $\eta=2/(\lambda_1+\lambda_n)$, el peor factor vale $(\kappa-1)/(\kappa+1)$, donde $\kappa=\lambda_n/\lambda_1$ es el \emph{número de condición} de $A$. Cuando $\kappa$ es grande, la tasa de aprendizaje está limitada por la dirección de mayor curvatura y el progreso en la de menor curvatura es muy lento. Es el fenómeno de mal condicionamiento descrito en~\cite[§4.2 y §4.3.1]{gbc}, y la motivación de métodos como Adam (módulo~7).
\end{observacion}

\begin{proposicion}[Paso óptimo a lo largo del gradiente]\label{prop:paso-optimo}
Sea $g=\nabla f(x)\ne 0$ y $H=H_f(x)$ con $g\T Hg>0$. La aproximación de segundo orden $q(\eta)=f(x)-\eta\,g\T g+\tfrac12\eta^{2}g\T Hg$ de $f(x-\eta g)$ es mínima en
\[
\eta^{*}=\frac{g\T g}{g\T Hg}.
\]
Si $f$ es cuadrática, $\eta^{*}$ minimiza exactamente $f(x-\eta g)$.
\end{proposicion}

\begin{proof}
$q$ es un polinomio de grado~2 en $\eta$ con coeficiente director $\tfrac12 g\T Hg>0$; su derivada $q'(\eta)=-g\T g+\eta\,g\T Hg$ se anula en $\eta^{*}$. Si $f$ es cuadrática, el resto del \cref{teo:taylor} es nulo y $q(\eta)=f(x-\eta g)$.
\end{proof}

\begin{ejemplo}
En el \cref{ej:dg}, $g=(2,6)\T$, $g\T g=40$ y $g\T Hg=2\cdot 4+6\cdot 36=224$, luego $\eta^{*}=40/224=5/28\approx 0{,}1786$. El paso lleva a $(9/14,\,-1/14)\approx(0{,}6429;\,-0{,}0714)$, con $f=3/7\approx 0{,}4286$, frente a $1{,}12$ con $\eta=0{,}1$.
\end{ejemplo}

\begin{observacion}[Relación con la física]
Cerca de un mínimo, $f(x^{*}+h)\approx f(x^{*})+\tfrac12h\T Hh$ es la energía potencial de un sistema de muelles con matriz de rigidez $H$. Los autovalores grandes corresponden a muelles rígidos, que oscilan deprisa; un paso de descenso demasiado largo en esa dirección sobrepasa el equilibrio y el sistema discreto diverge, como establece la \cref{prop:dg-cuadratica}.
\end{observacion}

\begin{lecturas}
\item Deisenroth, Faisal y Ong~\cite{mml}, §5.7 «Higher-Order Derivatives» y §5.8 «Linearization and Multivariate Taylor Series».
\item Goodfellow, Bengio y Courville~\cite{gbc}, §4.2 «Poor Conditioning» y §4.3.1 «Beyond the Gradient: Jacobian and Hessian Matrices».
\item MIT 18.02~\cite{mit1802}, clase~8 «Max-Min Problems» y clase~9 «2nd Derivative Test».
\item Dauphin \emph{et al.}~\cite{dauphin}, sobre la abundancia de puntos de silla en alta dimensión.
\end{lecturas}

\section*{Ejercicios del tema 4}
\addcontentsline{toc}{section}{Ejercicios del tema 4}

\begin{ejercicio}{1}{eje:4-vf}
Justifique si son verdaderas o falsas: (a) si $\nabla f(x)=0$, entonces $x$ es un mínimo local; (b) si $f\in C^{2}$, su hessiano es simétrico; (c) si $\nabla f(x^{*})=0$ y $H_f(x^{*})$ tiene un autovalor nulo y los demás positivos, $x^{*}$ es un mínimo local.
\end{ejercicio}

\begin{ejercicio}{2}{eje:4-taylor}
Sea $f(x,y)=x^{3}+xy$. En $(1,1)$ y con $h=(0{,}1;\,-0{,}2)$, calcule las aproximaciones de Taylor de primer y segundo orden de $f(x+h)$ y compárelas con el valor exacto. ¿Puede ser $f$ convexa en un entorno de $(1,1)$?
\end{ejercicio}

\begin{ejercicio}{2}{eje:4-mit}
Sea $f(x,y)=x^{3}-3x^{2}+xy-y^{2}+x+y$. Encuentre sus puntos críticos y clasifíquelos con el \cref{teo:criterio}.
\fuente{adaptado de MIT 18.02, primavera de 2006~\cite{mit1802}, \emph{Problem Set}~4, segunda parte, problema~1 (el original pide además los extremos en el cuadrado $[0,2]^{2}$).}
\end{ejercicio}

\begin{ejercicio}{2}{eje:4-tasa}
Sea $f(x,y)=x^{2}+4y^{2}$. (a) Determine las tasas de aprendizaje $\eta$ para las que el descenso de gradiente converge. (b) Desde $(1,1)$, calcule dos iteraciones y el valor de $f$ con $\eta=0{,}2$ y con $\eta=0{,}3$. (c) Calcule la tasa $\eta$ que minimiza el peor factor de contracción y el valor de ese factor.
\end{ejercicio}

\begin{ejercicio}{3}{eje:4-necesaria}
Sea $f\in C^{2}(U)$ y $x^{*}$ un mínimo local. Demuestre que $H_f(x^{*})\succeq 0$ (condición necesaria de segundo orden). Muestre con un ejemplo que el recíproco es falso: $\nabla f(x^{*})=0$ y $H_f(x^{*})\succeq 0$ no implican que $x^{*}$ sea un mínimo local.
\end{ejercicio}
